\chapter*{Preface}

Reduced-order models are usually introduced through a method: a basis, a projection, an error estimate. Using one on a real structure raises a different set of questions. Where do the matrices come from, and in what form? What does the reduced model need from the finite element code, and what does it have to be told? Which part of a run is paid once and which part again for every query? And when the reduced answer is wrong, how would you know? This book is about those questions. It describes how a reduced-order modelling package, \rom{}, is built on top of a finite element package, \fea{}, for the analysis of structures, and it checks, one claim at a time, that the two packages do what is said of them.

\section*{Who should read this book}

This book is not an introduction. We assume that you already know the finite element method well enough to write down the weak form of a beam or a plane-stress problem, to say what an element stiffness matrix is and how it is assembled, and to follow a Newton iteration. We assume that you have met reduced-order modelling: the idea of projecting a large system on a small basis, the proper orthogonal decomposition, modal truncation. And we assume that you read Python and NumPy without effort. Chapters on elements, dynamics, nonlinearity and contact therefore do not re-derive shape functions or time integrators. They concentrate on what is specific to the packages and on the places where a plausible result can be wrong.

The book is aimed at three kinds of reader. An engineer who has a finite element model and wants a fast surrogate of it will find which reduction suits which question, and what each costs offline and online. A developer who wants to build or extend a solver will find a structure in which a new element, load or backend is a subclass or a registry entry, with the tests that keep it honest. A researcher reproducing a published method will find nine such reproductions, including the ones that failed.

\section*{Scope of this book}

The subject is the practical implementation of reduced-order models over finite element analysis in structural mechanics. Here is what that includes. Full-order models of beams, plates, shells, plane and solid continua, in linear and nonlinear statics, in dynamics, and with frictionless unilateral contact. Reduced-order models built from these: POD--Galerkin with an affine parameter dependence, balanced truncation, Krylov and second-order Arnoldi reduction, frequency-domain models with greedy training and certified bounds, component mode synthesis, non-intrusive identification and nonlinear modal models, and hyper-reduction including contact. Both packages are used through plain arrays, and the reduced models never import the finite element code.

Several things are out of scope. The book is not a validation study: every comparison is against an exact solution, an independent computation or a published number, and none is against an experiment, so a small error shows that the code does what its equations say and not that the equations suit your structure. It does not compare the packages with commercial or open-source codes. It does not cover fluid mechanics, coupled multiphysics in general, friction, or the training of large neural networks. The optional PyTorch layer is described where it matters, and what needs it is said not to have been run when that is the case. Structured meshes are generated inside the package, and unstructured mesh generation is left to other tools.

The book is tied to one pinned release of the repository, \texttt{fea\_engine}~1.0.1 and \texttt{rom\_engine}~0.2.0 (\booktag), so that every listing, number and code anchor refers to the same code. Where the packages changed after the tag, the text says so.

\section*{Outline of the book}

The book has four parts and four appendices. Part~I, \emph{Architectural foundations}, is about structure. Chapter~\ref{ch:arch} shows the layers of \fea{}, the element interface through which solvers talk to elements, why \rom{} never imports \fea{}, and what the backend dispatcher costs. Chapter~\ref{ch:mesh} treats the meshes that elements sit on: topologies, quality checks, structured generators and grading, meshes of several element types, node sets, file formats, and the dispatcher that matches elements to meshes.

Part~II, \emph{Full-order physics}, builds the models that will be reduced. Chapter~\ref{ch:elem} goes from the weak form to element matrices and looks at locking, integration rules and the mass matrix. Chapter~\ref{ch:dyn} covers natural frequencies, implicit and explicit time integration, modal superposition and the frequency domain. Chapter~\ref{ch:nonlin} treats material and geometric nonlinearity and the methods that follow a response past a limit point. Chapter~\ref{ch:batch} is about making nonlinear assembly fast by batching. Chapter~\ref{ch:contact} enforces unilateral contact exactly with sparse active sets and a second-order time scheme.

Part~III, \emph{Reduced-order modelling}, is the centre of the book. Chapter~\ref{ch:pod} builds POD--Galerkin models with an affine parameter dependence and counts when the offline cost is repaid. Chapter~\ref{ch:bt} reduces a model through its input--output map by balanced truncation, Krylov projection and SOAR. Chapter~\ref{ch:freq} works in the frequency domain, chooses training frequencies by a greedy search, and tests a certified error bound. Chapter~\ref{ch:cms} reduces the parts of an assembly separately and joins them again. Chapter~\ref{ch:nonint} identifies models from data and fits only the nonlinear force of a structure that bends a lot. Chapter~\ref{ch:hyper} removes the last dependence on the full mesh by hyper-reduction, and carries contact into the reduced model.

Part~IV, \emph{Verification}, has one chapter. Chapter~\ref{ch:verif} says what kind of evidence each number in the book is, rebuilds nine published studies from the packages and reports how each came out, and ends with one complete pipeline from mesh to a prediction on loads not seen in training. The appendices cover differentiable assembly and automatic-differentiation tangents (Appendix~\ref{ch:diff}), the conventions and tolerances used in the book (Appendix~\ref{ch:conv}), a glossary (Appendix~\ref{ch:gloss}), and a map from chapters to the package tests that cover them (Appendix~\ref{ch:map}).

You do not have to read the book in order. Part~III depends on Part~II only for the models it reduces, so a reader who comes for the reduction methods can start at Chapter~\ref{ch:pod} and return to Chapters~\ref{ch:nonlin} and~\ref{ch:contact} when Chapters~\ref{ch:nonint} and~\ref{ch:hyper} need them. A reader who only wants to extend the finite element package can stop after Part~II.

\section*{Conventions used in this book}

Each chapter opens with a box that says what the chapter is about and what you should be able to do afterwards, and most open with a concrete problem before the theory. Short blue boxes inside the text are sidenotes: definitions and detours that can be skipped on a first reading. Listings headed \emph{Try it} are complete and runnable with the packages on the path, and the output printed under them is what the code produced. Algorithms are given as pseudocode, followed by a short paragraph on their cost, how they fail and where the code lives. Each chapter ends with a box on limits, which lists what was not tested, then a summary and exercises. A table at the end of each technical chapter maps its concepts to the code.

Code is cited as \anchor{importable.name}, and the last column of the code-map tables says whether the chapter's own checks call it. Numbers in the text are generated by scripts and enter the book as macros, so the number in the text is the one that was measured. They are single runs under one reference setup, and carry the caveats of the Reading guide that follows; the numbers of Chapter~\ref{ch:verif} on the nine reproductions are quoted from the author's reports and were not regenerated at the tag. The terms \emph{verification} and \emph{validation} are used in a fixed sense, set out in the section on how to read the verification sections.
